\section{Теоремы Ролля, Лагранжа и Коши.}

\subsection{Теорема Ролля}
\begin{theorem}[Теорема Ролля]
Пусть $y=f(x)$ и выполнены условия:
\begin{enumerate}
  \item $f$ непрерывна на $[a,b]$, обозначим \fbox{1}\;
  \item $f$ дифференцируема на $(a,b)$, обозначим \fbox{2}
  \item $f(a)=f(b)$.
\end{enumerate}
Тогда $\exists\, \xi \in (a,b)$ такое, что $f'(\xi)=0$.
\end{theorem}

\begin{Proof}
Пусть
\[
M=\sup_{[a,b]} f(x), \qquad m=\inf_{[a,b]} f(x).
\]

\begin{enumerate}
\item $M=m$. Тогда $f(x) =  \mathrm{const}$ на $[a,b]$, значит $f'(x)=0$ на $(a,b)$.
Возьмём, например,
\[
\xi=\frac{a+b}{2}.
\]

\item $M>m$. Кто-то из них достигается не в концах $a$ и $b$.
Пусть
\[
M=f(x_1), \qquad x_1\in(a,b).
\]
Тогда $x_1$ — точка максимума, и $\exists\, f'(x_1)$ (т.к. внутри функция дифф.)
используя теорему Ферма (необх. условие экстремума):
\[
f'(x_1)=0.
\]
\end{enumerate}
\end{Proof}


\subsection{Теорема Лагранжа}
\begin{theorem}[Теорема Лагранжа]
Если для $f$ выполнены \fbox{1}\ и \fbox{2}, то
\[
\exists\, \xi\in(a,b):\qquad \frac{f(b)-f(a)}{b-a}=f'(\xi).
\]
Графический смысл этой дроби - это угловой коэфицент. У нас найдётся точка, такая что у неё касательная будет параллельна вот этой так называемой секущей. 
\end{theorem}

\begin{Proof}
\[
F(x)=f(x)-\lambda x \quad \leftarrow \ f \text{ выполнены \fbox{1}\ и \fbox{2} на }[a,b] \text{ и } (a,b) \text{ соотвественно}
\]

\[
F(a)=F(b)\iff f(a)-\lambda a=f(b)-\lambda b
\iff \lambda=\frac{f(b)-f(a)}{b-a}
\]

По теореме Ролля для $y=F(x)$:
\[
\exists\,\xi\in(a,b):\quad F'(\xi)=0
= f'(\xi)-\lambda=0
\]

\[
f'(\xi)=\lambda
\]
\end{Proof}


\textbf{Следствия из теоремы Лагранжа.}

1) Если $y=f(x)$, $f$ выполнены \fbox{1}\ \ \fbox{2} и
$f'(x)=0$ на $(a,b)$, то $f(x)=\mathrm{const}$ на $[a,b]$.

\begin{Proof}
Предположим противное:
\[
\exists\,x_1,x_2\in[a,b]:\ f(x_1)\ne f(x_2).
\]
На $[x_1,x_2]$ выполнены \fbox{1}\ \ \fbox{2}, значит по теореме Лагранжа:
\[
\exists\,\xi\in(x_1,x_2):\qquad f(x_2)-f(x_1)=f'(\xi)(x_2-x_1).
\]
Но $f'(\xi)=0$, а левая $\ne 0$ — противоречие.
\end{Proof}

2) Если $f$ и $g$ \fbox{1}\ на  $[a,b]$, \fbox{2} на $(a,b)$ и
$f'(x)=g'(x)$ на $(a,b)$, то $f(x)-g(x)=\mathrm{const}$ на $[a,b]$.
\\
\\
Доказывается аналогично, нужно предъявить $h(x) = f(x) - g(x)$ и доказывать по первому следствию.

\subsection{Теорема Коши}
\begin{theorem}[Теорема Коши]
Если $f(x)$ и $g(x)$ выполнено \fbox{1} и \fbox{2} и $g(a)\ne g(b)$,
$g'(x)\ne 0$ на $(a,b)$, тогда:
\[
\exists\,\xi\in(a,b):\qquad
\frac{f(b)-f(a)}{g(b)-g(a)}=\frac{f'(\xi)}{g'(\xi)}.
\]
\end{theorem}

\begin{Proof}
Рассмотрим $F(x)=f(x)-\lambda g(x)$ — выполнено \fbox{1} и \fbox{2}.\\
Найдём $\lambda:\ F(a)=F(b)$
\[
\lambda=\frac{f(b)-f(a)}{g(b)-g(a)}.
\]

По теореме Ролля:
\[
\exists\,\xi\in(a,b):\quad F'(\xi)=0
\]
\[
\Downarrow
\]
\[
f'(\xi)-\lambda g'(\xi)=0
\]
\[
\frac{f'(\xi)}{g'(\xi)}=\lambda.
\]
\end{Proof}